% A conventional article: the common case.
\documentclass[11pt,a4paper]{article}
\usepackage[utf8]{inputenc}
\usepackage[T1]{fontenc}
\usepackage{amsmath,amssymb,amsthm}
\usepackage[margin=1in]{geometry}
\usepackage{graphicx}
\usepackage[colorlinks=true,linkcolor=blue]{hyperref}
\usepackage{natbib}

\title{On the Consistency of Estimators under Weak Dependence}
\author{A.\ Researcher\thanks{Department of Statistics.} \and B.\ Coauthor}
\date{\today}

\newtheorem{theorem}{Theorem}[section]
\newtheorem{lemma}[theorem]{Lemma}

\begin{document}
\maketitle

\begin{abstract}
We establish consistency of $\hat\theta_n$ under $\alpha$-mixing, relaxing the
i.i.d.\ assumption of \citet{knuth1984}. Simulations support the theory.
\end{abstract}

\section{Introduction}
\label{sec:intro}

Let $X_1,\dots,X_n$ be a sequence of random variables with common marginal
distribution $F$. Prior work \citep{knuth1984,lamport1994} assumes independence;
we do not. See Section~\ref{sec:results} for the main result.

\subsection{Notation}
\label{sec:notation}

Throughout, $\|\cdot\|$ denotes the Euclidean norm, and we write
\begin{equation}
  \label{eq:objective}
  Q_n(\theta) = \frac{1}{n} \sum_{i=1}^{n} q(X_i, \theta),
\end{equation}
with $\hat\theta_n = \arg\min_{\theta \in \Theta} Q_n(\theta)$.

\section{Results}
\label{sec:results}

\begin{theorem}[Consistency]
\label{thm:main}
Suppose $\Theta$ is compact and $q$ is continuous. Then
$\hat\theta_n \to \theta_0$ in probability.
\end{theorem}

\begin{proof}
The argument follows the standard uniform-convergence route. By
Lemma~\ref{lem:ulln},
\[
  \sup_{\theta \in \Theta} \left| Q_n(\theta) - Q(\theta) \right| \xrightarrow{p} 0 .
\]
Identification then gives the claim.
\end{proof}

\begin{lemma}[Uniform LLN]
\label{lem:ulln}
Under $\alpha$-mixing with $\sum_k \alpha(k) < \infty$, the empirical objective
converges uniformly.
\end{lemma}

\subsection{Discussion}

The conditions are mild. Note that:
\begin{itemize}
  \item compactness can be weakened to local compactness;
  \item continuity may be replaced by upper semicontinuity;
  \item the mixing rate is not sharp.
\end{itemize}

An enumerated version, for contrast:
\begin{enumerate}
  \item First, we verify identification.
  \item Second, we verify the uniform law.
  \item Third, we conclude.
\end{enumerate}

\begin{figure}[htbp]
  \centering
  \includegraphics[width=0.7\linewidth]{figures/convergence}
  \caption{Empirical convergence of $\hat\theta_n$ across 500 replications.}
  \label{fig:convergence}
\end{figure}

\bibliographystyle{plainnat}
\bibliography{references}

\end{document}
